% Appendices --- department template (v2.2 APPEND-01)

\section*{Appendix 1: Source Code and Project Artifacts}
\addcontentsline{toc}{section}{Appendix 1: Source Code and Project Artifacts}

The source code, training settings, and the files used to build this report are available at:

\begin{center}
\url{https://github.com/wikiepeidia/Internship-project}
\end{center}

The layout of the source tree and what each folder does is described in Chapter~IV. One point bears repeating: \texttt{src/model\_adaptation} holds the older experiment code. It is kept because the recorded results were produced with it, and it should not be read as the current design.

The exported Qwen model is not stored in the repository because it is a 4.28~GB file. A record kept beside it notes the file size, which checkpoint it came from, and the two occasions it was loaded to confirm that it works.

\section*{Appendix 2: Dataset Class Distribution}
\addcontentsline{toc}{section}{Appendix 2: Dataset Class Distribution}

Table~\ref{tab:dataset-stats} in Chapter~III gives the class counts for the three splits used in this work. What the quality checks found is reported in Chapter~V.

An earlier version of this work used a different 254-message validation set, shown in Table~\ref{tab:validation-class-dist} for historical interest only. Its numbers were never recalculated on the final dataset, and none of the results in this report come from it.

\begin{table}[H]
\centering
\caption{The older 254-message validation set, kept only as background.}
\label{tab:validation-class-dist}
\setlength{\tabcolsep}{10pt}
\renewcommand{\arraystretch}{1.3}
\begin{tabular}{lr}
\toprule
\textbf{Class} & \textbf{Messages} \\
\midrule
Bank impersonation       & 56  \\
Zalo social engineering  & 75  \\
Task scam                & 62  \\
Benign                   & 61  \\
\midrule
\textbf{Total}           & \textbf{254} \\
\bottomrule
\end{tabular}
\end{table}

\section*{Appendix 3: Where the Results Came From}
\addcontentsline{toc}{section}{Appendix 3: Where the Results Came From}

Table~\ref{tab:evaluation-snapshot} lists the main numbers in one place, with the run each of them came from.

\begin{table}[H]
\centering
\footnotesize
\renewcommand{\arraystretch}{1.12}
\begin{tabular}{L{3.4cm}L{5.2cm}L{4.7cm}}
\toprule
\textbf{What was measured} & \textbf{Result} & \textbf{What it does not say} \\
\midrule
The dataset & 2{,}097 messages; 1{,}658 / 219 / 220; no source article in two splits; largest article 7.9638\% & Nothing about how close it is to real scam traffic \\
Judge model & 1{,}395 of 2{,}097 passed (66.52\%); 1{,}561 scores reused, 536 scored fresh & Not a fresh independent scoring of the whole set \\
Manual review & 44 of 100 passed; agreed with the judge on 87 & A mixed sample chosen to cover every class and both judge outcomes, not a random one \\
Label corrections & 324 reviewed; 91 dropped; 233 agreed; 57 used & 176 agreed Zalo corrections were left out because they came from one article \\
Ordinary LoRA measurement & 31 steps; 53.274~s median; 7{,}902 of 8{,}151~MiB; no out-of-memory error & The 18.42~h figure is a projection from a part-run, not an observed time \\
QLoRA measurement & 40 measured steps; 3.462389~s median; 7{,}516 of 8{,}151~MiB; no out-of-memory error & Its 72-minute estimate is a projection, not an observed run \\
Qwen on validation & $N=219$; accuracy 0.990868; macro F1 0.988515 & Seed 42, checkpoint 200; used to select a checkpoint \\
PhoBERT on validation & $N=219$; accuracy 0.986301; macro F1 0.984893 & Seed 42, checkpoint 100; used to select a checkpoint \\
Qwen on held-out messages & $N=220$; 0.981818 / 0.980493 / 0.981848; 0 unreadable; 1 risky called safe & One run, run once \\
PhoBERT on held-out messages & $N=220$; 0.990909 / 0.990892 / 0.990925; 0 unreadable; 1 risky called safe & Higher here, but one run is not proof of a better model \\
Exported model & Q8\_0 GGUF, 4{,}280{,}403{,}232 bytes; loaded twice to confirm it works & PhoBERT was not exported; no deployment version was fitted \\
\bottomrule
\end{tabular}
\caption{All the main numbers in one place. The three figures for the held-out messages are accuracy, macro F1, and weighted F1. Measurements, projections, validation results, and final results are kept apart on purpose.}
\label{tab:evaluation-snapshot}
\end{table}


\section*{Appendix 4: What Went Wrong}
\addcontentsline{toc}{section}{Appendix 4: What Went Wrong}

Failures were kept in the record rather than removed. Table~\ref{tab:failure-recovery-ledger} lists what failed, what was done about it, and what the result may and may not be used to claim.

\small
\begin{longtable}{p{2.8cm}p{6.1cm}p{5.1cm}}
\caption{What failed, and what was done about it.}\label{tab:failure-recovery-ledger}\\
\toprule
\textbf{Stage} & \textbf{What went wrong} & \textbf{What was done, and what it can be used to claim} \\
\midrule
\endfirsthead
\toprule
\textbf{Stage} & \textbf{What went wrong} & \textbf{What was done, and what it can be used to claim} \\
\midrule
\endhead
Rebuilding the Zalo messages & A review found that one whole group of Zalo messages had been written as narration describing a scam instead of as a scam message. All 240 of them were affected. & They were rewritten from the same source articles, producing 300 replacements, of which four were rejected for drifting away from their source. The final set keeps 296. They are still generated messages, and they were written by the same model family as the judge that later scored them. \\
\addlinespace
Ordinary LoRA measurement & The first attempt stopped before the first step because a strict check expected an older library version. The corrected retry ran 31 steps and then ended on a driver fault. & It recorded no out-of-memory error, but memory reached 7{,}902 of 8{,}151~MiB with 9~MiB free, and it pointed at roughly eighteen hours of computation. It is evidence about the hardware only: it gives no accuracy figure and does not prove that LoRA is impossible. \\
\addlinespace
QLoRA measurement & The projection this run produced counted only one validation pass, not the 219-message generation that the full run repeats at every checkpoint, so it understated the real duration. & The measurement adapter was deleted and the full model was trained from a fresh start. The projection is used as evidence that the method fits the card, not as an observed run time. \\
\addlinespace
Full Qwen training & The first full attempt reached real checkpoints but exposed two mistakes: the output folder name left out the run ID, and one code path rebuilt the saved file through the wrong Python type. & It was stopped on purpose and kept. The accepted run restarted from step zero, completed 1{,}245 steps, selected step 200, and produced the exported model. \\
\addlinespace
Full PhoBERT training & PhoBERT had to stay an ordinary four-class classifier with nothing quantized, and any failed attempt had to be kept rather than quietly replaced. & The accepted run completed 312 steps and selected step 100. No PhoBERT GGUF was made. \\
\addlinespace
Starting the final evaluation & Five attempts to launch the final evaluation failed their own access checks before reading anything. & The sixth ran, read the 220 held-out messages once, evaluated both models, and then permanently disabled a repeat. Those five records describe those five attempts; they are not a full account of everything that ever opened the files. \\
\addlinespace
A correction about file access & An earlier draft said the held-out messages had been left untouched until the evaluation. That was wrong: automatic integrity tests had opened and checked the files before and after it. & The result itself was not changed or re-run. The defensible statement is that the models saw those messages exactly once, not that nothing had ever read the files. A written correction is kept with the result. \\
\addlinespace
Reorganising the code & The code was reorganised to be easier to read, but a later review recorded six critical and three warning findings. & It is reported as tidying work. It is not a claim that the system is secure, and it is not what produced the results in this report. \\
\bottomrule
\end{longtable}
\normalsize

\section*{Appendix 5: How to Reproduce This}
\addcontentsline{toc}{section}{Appendix 5: How to Reproduce This}

Every number in this report was written to a file by the run that produced it, and those files are in the repository. Table~\ref{tab:repro-evidence-map} says where to look for each group of claims.

\small
\begin{longtable}{p{3.0cm}p{6.3cm}p{4.6cm}}
\caption{Where each group of claims comes from.}\label{tab:repro-evidence-map}\\
\toprule
\textbf{Claim} & \textbf{Where it comes from} & \textbf{What it may be used for} \\
\midrule
\endfirsthead
\toprule
\textbf{Claim} & \textbf{Where it comes from} & \textbf{What it may be used for} \\
\midrule
\endhead
The dataset and its quality & The dataset manifest \path{data/manifests/manifest.json} and the judge summary \path{data/processed/judge-summary.json}, together with the record of the manual review. & The class counts and the quality scores. The manual review covers a sample, not the whole set. \\
\addlinespace
Choosing between LoRA and QLoRA & The two measurement-run reports, which record step times and memory. & Time and memory only. Neither run produced a trained model. \\
\addlinespace
Validation comparison & The validation comparison record, covering the same 219 messages for both models. & Comparing quality. It does not support comparing speed. \\
\addlinespace
The exported model & The export record kept beside the GGUF file. & That the Q8\_0 file was produced from the selected checkpoint and loads. There is no PhoBERT GGUF. \\
\addlinespace
The final evaluation & The evaluation record for the 220 held-out messages. & The final numbers for both models. The evaluation was run once and cannot be repeated. \\
\addlinespace
The correction about file access & The written correction kept beside the final result. & It belongs with that result and should be read alongside it. \\
\addlinespace
How the code is organised & \path{architecture/module-boundaries.json} and \path{docs/architecture/overview.md}. & The current layout and the installed commands. It is not what produced the results. \\
\addlinespace
Open problems in the code & The latest code review. & Six critical and three warning findings, still open. \\
\bottomrule
\end{longtable}
\normalsize
